\documentclass{article}

\usepackage[margin=1in]{geometry}
\usepackage{amsmath,amsthm,amssymb}
\usepackage[spanish]{babel}
\usepackage{enumitem}
\usepackage{fancyhdr}
\usepackage{lastpage}

% Number sets
\newcommand{\R}{\mathbb{R}}
\newcommand{\Z}{\mathbb{Z}}
\newcommand{\N}{\mathbb{N}}
\newcommand{\Q}{\mathbb{Q}}
\newcommand{\C}{\mathbb{C}}
\newcommand{\K}{\mathbb{K}}

% Functional analysis operators
\newcommand{\norm}[1]{\left\lVert #1 \right\rVert}       % Norm

% Exercise environment
\newenvironment{ejercicio}[1]{\begin{trivlist}
\item[\hskip \labelsep {\bfseries Ejercicio}\hskip \labelsep {\bfseries #1.}]}{\end{trivlist}}

\setlist[enumerate,1]{label=\alph*), leftmargin=*, itemsep=2pt}
\setlist[enumerate,2]{label=\roman*., leftmargin=*, itemsep=2pt}

\pagestyle{fancy}
\fancyhf{}
\fancyhead[L]{Departamento de Matemática - Facultad de Ciencias Exactas - UNLP}
\fancyfoot[L]{Análisis funcional 2026}
\fancyfoot[R]{Página \thepage\ de \pageref{LastPage}}

\begin{document}

\begin{center}
    {\Huge Análisis funcional}\\[0.15cm]
    Práctica de repaso

    \vspace{0.4cm}

    {\Large Espacios vectoriales de dimensión finita}\\[0.1cm]
    \textit{Recordando un poco de Álgebra Lineal}
\end{center}

\noindent \textbf{Notación:} Durante en toda la práctica vamos a considerar a \(\mathcal{V}\) un \(\K\)-espacio vectorial de dimensión finita. Denotaremos \(L(\mathcal{V})\) al espacio vectorial de operadores lineales \(A:\mathcal{V}\to\mathcal{V}\). Dado un \(A\in L(\mathcal{V})\), un número complejo \(\lambda\in\C\) se dice autovalor de \(A\) si existe \(x\in\mathcal{V}\) no nulo tal que \(Ax=\lambda x\). Al conjunto de autovalores lo denominaremos espectro y lo simbolizaremos \(\sigma(A)\). Y por medio de \(A^*\) se denotará al operador adjunto de \(A\).

\begin{ejercicio}{1}
    Dado \(A\in L(\mathcal{V})\), probar lo siguiente:
    \begin{enumerate}
        \item \(\displaystyle \norm{A} = \sup_{\substack{x\in\mathcal{V}\\x\neq0}} \frac{\norm{Ax}}{\norm{x}}\) es una norma en \(L(\mathcal{V})\).
        \item Probar que \(\displaystyle \norm{A} = \sup_{\substack{x\in\mathcal{V}\\\norm{x}=1}} \norm{Ax} = \text{\'inf}\{c\in\R^+ : \norm{Ax}\leq c\norm{x}\}\).
    \end{enumerate}
\end{ejercicio}

\vspace{0.5cm}

\begin{ejercicio}{2}
    Sea \((\mathcal{V}, \langle\cdot,\cdot\rangle)\) un espacio con producto interno, y sea \(\norm{\cdot}\) la norma que induce dicho producto, dada por
    \[
        \norm{x} = \langle x,x\rangle^{1/2}.
    \]
    Demostrar que, para cualquier \(A\in L(\mathcal{V})\), se verifica lo siguiente:
    \begin{enumerate}
        \item \(\displaystyle \norm{A} = \sup_{\substack{x,y\in\mathcal{V}\\\norm{x}=1,\norm{y}=1}} |\langle Ax,y\rangle|\).
        \item \(\norm{A} = \norm{A^*}\).
        \item Si \(A=A^*\) entonces \(\displaystyle \norm{A} = \sup_{\substack{x\in\mathcal{V}\\\norm{x}=1}} |\langle Ax,x\rangle| = \text{m\'ax}\{|\lambda| : \lambda\in\sigma(A)\}\).
        \item \(\norm{A}^2 = \norm{A^*A} = \text{m\'ax}\{|\lambda| : \lambda\in\sigma(A^*A)\}\).
    \end{enumerate}
\end{ejercicio}

\vspace{0.5cm}

\begin{ejercicio}{3}
    \begin{enumerate}
        \item Considerar la matriz
              \[
                  M = \frac{1}{5}\begin{pmatrix} 1 & 2 \\ 2 & 4 \end{pmatrix}
              \]
              Probar que \(M\) es la matriz asociada a la proyección lineal de \(\R^2\) en \(\R^2\) en la recta \(y=2x\).
        \item Probar que esta matriz es una proyección ortogonal.
        \item Sea \(T:\R^2\to\R^2\) dada por \(T(x,y) = (x,0)\). Encontrar la matriz \(P_1\) asociada a esta transformación lineal y probar que \(P_1 = P_1^2 = P_1^*\).
        \item Sea
              \[
                  P_2 = \begin{pmatrix} 0 & 0 \\ 3 & 1 \end{pmatrix}
              \]
              Probar que \(P_2 = P_2^2\) pero \(P_2 \neq P_2^*\). Encontrar la transformación lineal asociada a la matriz \(P_2\).
    \end{enumerate}
\end{ejercicio}

\vspace{0.5cm}

\noindent \textbf{Nota:} A partir de acá vamos a considerar \((\mathcal{V},\langle\cdot,\cdot\rangle)\) un espacio vectorial con producto interno y \(A\in L(\mathcal{V})\).

\begin{ejercicio}{4}
    Demostrar que \(R(A^*) = {(\ker A)}^\perp\) y que \(\ker A^* = {(R(A))}^\perp\).
\end{ejercicio}

\vspace{0.5cm}

\begin{ejercicio}{5}
    Probar lo siguiente:
    \begin{enumerate}
        \item \(\ker(A^*A) = \ker(A)\).
        \item Si \(A\) es normal entonces \(\ker(A^*) = \ker(A)\).
        \item Si \(A\) es normal y nilpotente entonces \(A=0\).
        \item Si \(A\) es normal y \(\sigma(A) = \{\lambda\}\) entonces \(A=\lambda I\).
    \end{enumerate}
\end{ejercicio}

\vspace{0.5cm}

\begin{ejercicio}{6}
    Sea \(\mathcal{S}\) un subespacio de \(\mathcal{V}\). Demostrar que si \(\mathcal{S}\) es invariante por \(A\) (\(A(\mathcal{S})\subset\mathcal{S}\)) entonces \(\mathcal{S}^\perp\) es invariante por \(A^*\).
\end{ejercicio}

\vspace{0.5cm}

\begin{ejercicio}{7}
    Sea \(\norm{\cdot}\) la norma sobre \(\mathcal{V}\) inducida por el producto interno. Demostrar que \(U\in L(\mathcal{V})\) es unitario si y sólo si \(\norm{Ux} = \norm{x}\) para todo \(x\in\mathcal{V}\).
\end{ejercicio}

\vspace{0.5cm}

\begin{ejercicio}{8}
    Probar que:
    \begin{enumerate}
        \item Si \(A\) es autoadjunto, entonces \(\sigma(A)\subseteq\R\). ¿Qué pasa si \(A\) es positivo?
        \item Si \(A\) es unitario, entonces \(\sigma(A)\subseteq S^1\).
    \end{enumerate}
\end{ejercicio}

\vspace{0.5cm}

\begin{ejercicio}{9}
    Probar que las siguientes afirmaciones son equivalentes:
    \begin{enumerate}
        \item \(A\geq0\).
        \item \(A=A^*\) y \(\sigma(A)\subseteq\R^+\).
        \item Existe \(B\geq0\) tal que \(A=B^2\).
        \item Existe \(C\in L(\mathcal{V})\) tal que \(A=C^*C\).
    \end{enumerate}
\end{ejercicio}

\vspace{0.5cm}

\begin{ejercicio}{10}
    Sean \(A,B\in L(\mathcal{V})\). Demostrar que \(\sigma(AB)\setminus\{0\} = \sigma(BA)\setminus\{0\}\).
\end{ejercicio}

\end{document}
